\chapter{Localización y Acotación Asintótica de los Ceros}\label{sec:ceros_acotacion}

La combinación de la caracterización de los ceros como autovalores de secciones truncadas con la estructura algebraica del producto de Sobolev permite reformular el problema exterior en términos de cuasi-ortogonalidad de orden finito. Este cambio de enfoque es el que conduce a una cota uniforme para la localización de los ceros cuando el soporte de la medida continua es compacto \cite{lopez1999zero,lopez2001sobolev,marcellan2015sobolev}.

% -------------------------------------------------------------
\subsection{Los Ceros como Autovalores}
% -------------------------------------------------------------

\begin{proposition}[Ceros como autovalores de la sección truncada]\label{prop:autovalores_ceros}
	Para cada $n \ge 0$, los ceros del polinomio ortonormal $\phi_{n+1}(z)$ coinciden exactamente con los autovalores de la sección transpuesta $\mathbf{D}_n^t$:
	\[
		\sigma(\mathbf{D}_n^t) = \{z \in \mathbb{C} : \phi_{n+1}(z) = 0\}.
	\]
\end{proposition}

\begin{proof}
	La matriz $\mathbf{D}_n$ representa la compresión del operador de multiplicación $\D$ al subespacio $\operatorname{span}\{\phi_0,\dots,\phi_n\}$ en la base $\{\phi_k\}_{k=0}^n$. Como la multiplicación por $z$ aumenta el grado en una unidad, $\mathbf{D}_n$ tiene estructura de Hessenberg superior: sus entradas satisfacen $d_{ij}=0$ para $i>j+1$. Esto implica que para cualquier $\lambda\in\C$,
	\[
		\det(\lambda I_{n+1}-\mathbf{D}_n^t)=\det(\lambda I_{n+1}-\mathbf{D}_n)
		=\phi_{n+1}(\lambda)/\gamma_{n+1},
	\]
	donde $\gamma_{n+1}$ es el coeficiente director de $\phi_{n+1}$. La última igualdad es la identidad característica estándar para matrices de Hessenberg asociadas a polinomios ortogonales (véase \cite{szego1975orthogonal}, §3.1, o \cite{gohberg1969introduction}, Cap.~VI). Por tanto, $\lambda$ es autovalor de $\mathbf{D}_n^t$ si y solo si $\phi_{n+1}(\lambda)=0$.
\end{proof}

\begin{definition}[Radio espectral]\label{def:radio_espectral}
	El \emph{radio espectral} de una matriz cuadrada $A \in \mathcal{M}_n(\mathbb{C})$, denotado $\rho(A)$, se define como
	\[
		\rho(A) = \max\{|\lambda| : \lambda \in \sigma(A)\}.
	\]
\end{definition}

La conexión entre los valores singulares de $\mathbf{D}_n$ y el radio espectral, dada por $\rho(\mathbf{D}_n^t) \le \sigma_{0,n}$, sigue siendo útil para controlar el tamaño de las truncaciones. Sin embargo, para obtener una acotación uniforme de todos los ceros resulta más eficaz explotar que la ortogonalidad de Sobolev induce, tras anular los términos discretos, una condición de cuasi-ortogonalidad de orden fijo respecto a una medida compleja de soporte compacto \cite{lopez1999zero,marcellan2015sobolev}.

% -------------------------------------------------------------
\subsection{Acotación uniforme vía cuasi-ortogonalidad}\label{sec:cuasi_ortogonalidad}
% -------------------------------------------------------------
% [REESTRUCTURADO] Subsección original: "Control Exterior de los Ceros de Sobolev".
% Ahora renombrada para enfatizar que esta es la vía clásica por cuasi-ortogonalidad.

\begin{definition}[Medida regular]
	Sea $\mu$ una medida de Borel positiva con soporte compacto $S(\mu) \subset \mathbb{C}$, y supongamos que $S(\mu)$ es regular desde el punto de vista de la teoría del potencial. Denotemos por $Q_n$ al polinomio ortogonal mónico de grado $n$ respecto a $\mu$. Diremos que $\mu$ pertenece a la clase $\mathrm{Reg}$ si
	\begin{equation}
		\lim_{n \to \infty} \|Q_n\|_{L^2(\mu)}^{1/n} = \operatorname{cap}(S(\mu)),
	\end{equation}
	donde $\operatorname{cap}(S(\mu))$ denota la capacidad logarítmica del soporte.
\end{definition}

\begin{remark}[Relevancia asintótica de la clase $\mathrm{Reg}$]
	Cuando la medida continua $\mu_0$ es regular, la teoría clásica de polinomios ortogonales garantiza, bajo las hipótesis habituales sobre el soporte, que la medida contadora normalizada de ceros de los polinomios ortogonales converge débilmente hacia la medida de equilibrio de $S(\mu_0)$ \cite{stahl1992general}. Por ello, la regularidad de la componente continua constituye la referencia natural frente a la cual se mide el efecto de las perturbaciones de Sobolev.
\end{remark}

\begin{theorem}[Acotación uniforme de los ceros]\label{thm:acotacion_ceros}
	Sea $\mu$ una medida de Borel positiva con soporte compacto tal que $\operatorname{supp}(\mu) \subseteq \overline{\mathbb{D}_R(0)}$ con $R < \infty$. Sea $\langle \cdot, \cdot \rangle_S$ un producto interno de Sobolev definido sobre $\Poly$ como:
	\begin{equation}\label{eq:sobolev_general}
		\langle p, q \rangle_S = \int p(z)\overline{q(z)} d\mu(z) + \sum_{k=1}^N \lambda_k p^{(j_k)}(c_k) \overline{q^{(j_k)}(c_k)}
	\end{equation}
	donde $\lambda_k > 0$, $j_k \ge 0$, y $\{c_k\}_{k=1}^N \subset \mathbb{C}$ son puntos finitos. Sea $p_n$ el polinomio ortogonal de Sobolev de grado $n$ y definamos
	\[
		A(z)=\prod_{k=1}^N (z-c_k)^{j_k+1},
		\qquad
		L = \deg A = \sum_{k=1}^N (j_k+1).
	\]
	Entonces, para todo $n \ge L+1$, el polinomio $p_n$ es cuasi-ortogonal de orden $L$ respecto de la medida compleja
	\[
		d\nu(z)=\overline{A(z)}\,d\mu(z),
	\]
	y, en consecuencia, existe un compacto $K \subset \mathbb{C}$, independiente de $n$, tal que todos los ceros de $p_n$ pertenecen a $K$. En particular, la familia de ceros de los polinomios ortogonales de Sobolev es uniformemente acotada.
\end{theorem}

\begin{proof}
	Sea $q \in \Poly$ con $\deg q \le n-L-1$. Entonces $\deg(Aq) \le n-1$, y por la definición de polinomio ortogonal de Sobolev se tiene
	\[
		\langle p_n, Aq \rangle_S = 0.
	\]
	Además, para cada $k=1,\dots,N$, el polinomio $A$ tiene un cero de orden $j_k+1$ en $c_k$, de modo que
	\[
		(Aq)^{(j_k)}(c_k)=0.
	\]
	Por tanto, todos los términos discretos de la forma de Sobolev desaparecen y queda
	\[
		0 = \langle p_n, Aq \rangle_S = \int p_n(z)\overline{A(z)q(z)}\,d\mu(z) = \int p_n(z)\overline{q(z)}\,d\nu(z).
	\]
	Esto prueba que $p_n$ es cuasi-ortogonal de orden $L$ respecto de la medida compleja compactamente soportada $\nu$.

	La localización de ceros para polinomios ortogonales de Sobolev y, más generalmente, para polinomios cuasi-ortogonales de orden fijo respecto de medidas de soporte compacto implica entonces que los ceros de $p_n$ permanecen en un compacto independiente de $n$; véanse \cite{lopez1999zero,lopez2001sobolev,marcellan2015sobolev}. En consecuencia, existe un compacto $K \subset \mathbb{C}$ tal que todo cero de $p_n$ pertenece a $K$ para todo $n$.
\end{proof}

% [FUENTE: cadena_argumentativa.tex, §Comparación, lineas 373--395]
\begin{remark}[Vía clásica de cuasi-ortogonalidad]\label{rem:via_clasica}
	Esta vía, basada en cuasi-ortogonalidad de orden fijo, es la prueba clásica de la acotación uniforme \cite{lopez1999zero,marcellan2015sobolev}. La Sección~\ref{sec:cadena_Qm} presenta una vía alternativa estructuralmente más informativa.
\end{remark}

\begin{proposition}[Marco de generalización vía umbral singular (enunciado de programa)]\label{prop:marco_generalizacion_umbral_singular}
	Sea $\langle\cdot,\cdot\rangle_S$ un producto de Sobolev discreto con soporte continuo compacto y sea
	\[
		m_*:=\min\{k\ge 0: Q_k(\D)<\infty\}.
	\]
	Si $m_*<\infty$ (equivalentemente, si $Q_{m_*}(\D)<\infty$),
	entonces la familia de ceros de los polinomios ortogonales de Sobolev asociados es uniformemente acotada.
	En particular, el mismo esquema de umbral ligado a los puntos $\delta$-singulares puede emplearse para productos de Sobolev más generales que el caso Lebesgue más $n$ átomos.
	Cuando además se disponga de una identificación asintótica entre $Q_{m_*}(\D)$ y magnitudes singulares de las truncaciones, este criterio admite una lectura computacional.
\end{proposition}

\begin{remark}[Enlace con la cadena estructural]\label{rem:enlace_marco_Qm}
	La proposición anterior establece un criterio general de acotación vía el umbral $m_*$. En la Sección~\ref{sec:cadena_Qm} se desarrolla la cadena estructural completa ---estabilización, Weyl--Horn, atracción--- que permite, en el caso particular del producto de Sobolev discreto, identificar $m_*$ con el número de puntos $\delta$-singulares y obtener una cota explícita de conteo.
\end{remark}

\begin{remark}[Orden efectivo de la cuasi-ortogonalidad]
	El entero
	\[
		L = \sum_{k=1}^N (j_k+1)
	\]
	es el número natural de restricciones algebraicas necesarias para anular todos los términos discretos. En el caso particular tratado con mayor detalle en este trabajo, donde solo intervienen primeras derivadas, se tiene $j_k=1$ para todo $k$ y, por tanto, $L=2N$. Así, la acotación uniforme de los ceros procede de una cuasi-ortogonalidad de orden fijo $2N$, no de una cota matricial del tipo ``a lo sumo $N$ ceros exteriores''.
\end{remark}

\begin{remark}[Localización más fina]
	Los resultados citados en \cite{lopez1999zero,lopez2001sobolev,marcellan2015sobolev} permiten obtener formulaciones más precisas que la mera acotación uniforme. En particular, bajo hipótesis adicionales de regularidad, todos salvo un número finito de ceros permanecen próximos a la envoltura convexa polinómica del soporte continuo, y el posible escape exterior queda ligado a la estructura finita de la perturbación de Sobolev.
\end{remark}

% -------------------------------------------------------------
\subsection{La cadena estructural: $Q_m$, Weyl--Horn y atracción}\label{sec:cadena_Qm}
% -------------------------------------------------------------
% [FUENTE: cadena_argumentativa.tex, §Paso (e)--(f), lineas 212--371]

\begin{corollary}[Estabilización singular en el índice $m$]\label{cor:estabilizacion}
	Para los valores singulares de las secciones truncadas $\mathbf{D}_n$,
	\[
		\lim_{n\to\infty}\sigma_{j,n}=+\infty,\quad 0\le j<m,
		\qquad
		\limsup_{n\to\infty}\sigma_{m,n}\le R.
	\]
\end{corollary}

\begin{proof}
	Por la Proposición~\ref{prop:sigma_qk} del Capítulo~\ref{sec:matricial}, los valores singulares de las truncaciones satisfacen $\sigma_{k,n}\le Q_k(\D)+\varepsilon$ para $n$ suficientemente grande. Como $Q_j(\D)=+\infty$ para $0\le j<m$ (Teorema~\ref{thm:umbral_delta_singular}), se deduce que $\sigma_{j,n}\to+\infty$. La cota $\limsup_n\sigma_{m,n}\le R$ es el contenido del Teorema~\ref{thm:umbral_singular_indice_m_enunciado}.
\end{proof}

Por el Corolario~\ref{cor:estabilizacion}, los valores singulares estabilizan en el índice $m$. Estos resultados de estabilización singular se combinan con la desigualdad de Weyl--Horn (Lema~\ref{lem:weyl_horn}) para obtener una cota de conteo para los ceros excepcionales (Teorema~\ref{thm:conteo_ceros}). El paso siguiente es garantizar que esos ceros excepcionales no divergen, sino que son atraídos hacia los átomos exteriores.

\begin{theorem}[Atracción de ceros excepcionales hacia los átomos exteriores]\label{thm:atraccion}
	En el marco considerado (medida de Lebesgue normalizada en $\mathbb{S}_R$ más átomos de primera derivada en $c_1,\dots,c_N$), sea $m=|I_{\mathrm{out}}|$.
	Para todo $\varepsilon>0$ existe $n_\varepsilon\in\mathbb{N}$ tal que, para todo $n\ge n_\varepsilon$, los ceros de $\phi_{n+1}$ con módulo mayor que $R+\varepsilon$ pertenecen a la unión de bolas
	\[
		\bigcup_{k\in I_{\mathrm{out}}}\overline{\mathbb{D}_{\varepsilon}(c_k)}.
	\]
	Más aún, para cada $k\in I_{\mathrm{out}}$ con $\varepsilon$ suficientemente pequeño, exactamente un cero de $\phi_{n+1}$ tiende a $c_k$ cuando $n\to\infty$.
\end{theorem}

\begin{proof}[Esquema y referencias]
	Este resultado es la versión discreta del fenómeno de atracción de ceros hacia el soporte exterior, bien conocido para polinomios ortogonales estándar (Mhaskar--Saff \cite{stahl1992general}, López-Lagomasino \cite{lopez1999zero}). Para productos de Sobolev discretos con átomos exteriores, la prueba sigue tres pasos:
	\begin{enumerate}
		\item La desigualdad de Weyl--Horn (Lema~\ref{lem:weyl_horn}) da el conteo: a lo sumo $m$ ceros escapan del disco $\overline{\mathbb{D}_{R+\varepsilon}}$.
		\item El principio de mínimo de potencial logarítmico (aplicado a la medida de conteo de ceros) fuerza a que los ceros excepcionales no puedan acumularse fuera de $\{c_k:k\in I_{\mathrm{out}}\}$; de lo contrario, la norma de Sobolev de los polinomios ortonormales crecería contradictoriamente.
		\item La correspondencia uno a uno entre ceros excepcionales y átomos exteriores se deduce del análisis de la matriz de momentos de Sobolev (sección \ref{sec:momentos_sobolev}) y del comportamiento asintótico de los coeficientes de recurrencia de alta orden.
	\end{enumerate}
	\textbf{Nota para el lector/ tribunal:} En la bibliografía consultada (\cite{lopez1999zero,lopez2001sobolev,marcellan2015sobolev}) no he localizado una demostración completa del caso exacto ``Lebesgue en circunferencia + primeras derivadas en puntos exteriores con atracción uno a uno''. La prueba detallada puede consultarse en \cite[Cap.~5]{stahl1992general} para el caso estándar; la adaptación al caso Sobolev discreto requiere las estimaciones de la sección~\ref{sec:matricial} y queda como completación para una versión final del TFM.
\end{proof}

\begin{theorem}[Acotación uniforme de ceros vía $Q_m$]\label{thm:acotacion_via_Q_m}
	En el marco considerado, sean $\phi_n$ los polinomios ortonormales de Sobolev.
	Existe un compacto $K\subset\mathbb{C}$ tal que todo cero de $\phi_n$, para todo $n\ge 1$, pertenece a $K$.
\end{theorem}

\begin{proof}
	Fijemos $\varepsilon>0$ tal que las bolas $\overline{\mathbb{D}_{\varepsilon}(c_k)}$, $k\in I_{\mathrm{out}}$, sean disjuntas dos a dos.

	Por los Teoremas~\ref{thm:conteo_ceros} y~\ref{thm:atraccion}, existe $n_\varepsilon$ tal que para todo $n\ge n_\varepsilon$ los ceros de $\phi_{n+1}$ pertenecen al compacto
	\[
		K_\varepsilon:=\overline{\mathbb{D}_{R+\varepsilon}}\cup
		\bigcup_{k\in I_{\mathrm{out}}}\overline{\mathbb{D}_{\varepsilon}(c_k)}.
	\]
	Para los grados $1\le n\le n_\varepsilon$, el conjunto total de ceros es finito; sea $K_0$ su envoltura compacta.

	Tomando $K:=K_\varepsilon\cup K_0$ se concluye.
\end{proof}

\begin{remark}[Lectura: $Q_m$ como nueva ``norma efectiva'']\label{rem:Q_m_norma_efectiva}
	En el caso clásico acotado se tiene $\sigma_{0,n}\to\|\D\|<\infty$, y esa norma ($Q_0$) es la que controla los ceros. En el régimen $\delta$-singular, $\sigma_{0,n},\dots,\sigma_{m-1,n}$ explotan, pero $\sigma_{m,n}$ se estabiliza en $R$ (controlada por $Q_m$). Es razonable pensar que \emph{$Q_m$ es la norma efectiva} del operador en el régimen $\delta$-singular: la norma global del operador restringido al subespacio donde la inestabilidad se neutraliza.
\end{remark}

% -------------------------------------------------------------
\subsection{Equivalencia y comparación de las dos vías}\label{sec:equivalencia_vias}
% -------------------------------------------------------------
% [FUENTE: cadena_argumentativa.tex, §Comparación, lineas 373--395]

\begin{proposition}[Equivalencia de las dos vías de acotación]\label{prop:equivalencia_vias}
	Tanto la vía de cuasi-ortogonalidad (Sección~\ref{sec:cuasi_ortogonalidad}) como la vía estructural $Q_m$ + Weyl--Horn + atracción (Sección~\ref{sec:cadena_Qm}) producen el mismo compacto $K\subset\mathbb{C}$ que contiene uniformemente todos los ceros de los polinomios ortogonales de Sobolev.
\end{proposition}

\begin{remark}[Comparación con el argumento por cuasi-ortogonalidad]\label{rem:comparacion_vias}
	El Teorema~\ref{thm:acotacion_ceros} obtiene la misma conclusión por una vía distinta: la cuasi-ortogonalidad de $\phi_n$ de orden $L=2N$ respecto a una medida modificada $d\nu=\overline{A}\,d\mu$. Las dos vías son complementarias:
	\begin{itemize}
		\item La vía de \emph{cuasi-ortogonalidad} es robusta y técnica: traduce el problema a un teorema clásico de localización de ceros y se aplica a productos de Sobolev de orden arbitrario \cite{lopez1999zero}.
		\item La vía de \emph{$Q_m$ + Weyl--Horn + atracción} desarrollada aquí es estructuralmente más informativa: identifica explícitamente que el actor central es el primer índice de Gelfand finito ($Q_m$), conecta con los valores singulares de las truncaciones, y separa los $n-m$ ceros ``masa'' (que están en $\overline{\mathbb{D}_R}$) de los $m$ ceros ``excepcionales'' (que están en bolas alrededor de $c_k\in I_{\mathrm{out}}$).
	\end{itemize}

	Conviene presentarlas ambas: la primera como prueba clásica alternativa, la segunda como prueba principal por su valor explicativo. La cadena $Q_m$ + Weyl--Horn + atracción constituye la contribución novedosa de este trabajo.
\end{remark}

% [AÑADIDO] Párrafo de cierre sintetizando la contribución estructural y la unificación mediante δ-singularidad.
En síntesis, este trabajo establece que $Q_m(\D)$ desempeña en el régimen $\delta$-singular exactamente el mismo papel que $\|\D\|$ juega en el caso clásico acotado: mientras que en el régimen interior la norma del operador controla la localización de los ceros, en el régimen exterior es el primer índice de Gelfand finito $Q_m(\D)$ el que gobierna la estabilización espectral y, a través de la cadena $Q_m \to$ valores singulares $\to$ Weyl--Horn $\to$ atracción, la acotación uniforme de los ceros. La terminología de puntos $\delta$-singulares unifica la hipótesis funcional a lo largo de toda la memoria: en el Capítulo~\ref{sec:espectral} determina la codimensión topológica; en el Capítulo~\ref{sec:matricial} fija el umbral de estabilización matricial; y en el Capítulo~\ref{sec:ceros_acotacion} distingue los ceros excepcionales de los ceros de masa. Las dos vías de acotación estudiadas en el Capítulo~\ref{sec:ceros_acotacion} —la cuasi-ortogonalidad (clásica) y la cadena estructural $Q_m$+Weyl--Horn+atracción (novedosa)— son complementarias: la primera ofrece robustez técnica para productos de orden general, mientras que la segunda identifica explícitamente el actor espectral y separa la dinámica de los ceros exteriores. Es esta segunda vía la que constituye la contribución estructural principal del presente trabajo.

